\documentclass[../../main.tex]{subfiles}
\begin{document}


{\bf ［解］}

長針，短針の長さを $R, r$ とする．

\begin{center}
\begin{tikzpicture}[scale=0.8, >=stealth]
  % 2時
  \begin{scope}[shift={(0,0)}]
    \draw[thick] (0,0) circle (1);
    \draw[thick] (0,0) circle (0.5);
    \draw[->, thick] (0,0) -- (0, 0.9);
    \draw[->, thick] (0,0) -- (60:0.45);
    \node at (0, -1.4) {2時(成す角$\frac{\pi}{3}$)};
  \end{scope}
  % 2時半
  \begin{scope}[shift={(5,0)}]
    \draw[thick] (0,0) circle (1);
    \draw[thick] (0,0) circle (0.5);
    \draw[->, thick] (0,0) -- (0, -0.9);
    \draw[->, thick] (0,0) -- (75:0.45);
    \node at (0, -1.4) {2時半(成す角$\frac{7}{12}\pi$)};
  \end{scope}
  % 4時
  \begin{scope}[shift={(10,0)}]
    \draw[thick] (0,0) circle (1);
    \draw[thick] (0,0) circle (0.5);
    \draw[->, thick] (0,0) -- (0, 0.9);
    \draw[->, thick] (0,0) -- (-30:0.45);
    \node at (0, -1.4) {4時(成す角$\frac{2}{3}\pi$)};
  \end{scope}
\end{tikzpicture}
\end{center}

題意から，2時，2時半の状態に余弦定理を用いると，
\begin{align}
\begin{cases}
4^2 = R^2 + r^2 - 2Rr \cos \frac{\pi}{3} = R^2 + r^2 - Rr  \\
6^2 = R^2 + r^2 - 2Rr \cos \frac{7\pi}{12}
\end{cases} \label{eq:1}
\end{align}
である．

ここで$\cos \frac{7}{12}\pi < 0$ だから半角公式より
\begin{align}
\cos \frac{7}{12}\pi 
  &= -\sqrt{\frac{1 + \cos \frac{7\pi}{6}}{2}} \\
  &= -\sqrt{\frac{1 - \frac{\sqrt{3}}{2}}{2}} \\
  &= -\frac{\sqrt{3}-1}{2\sqrt{2}}
 \end{align}
だから，\cref{eq:1}に代入して
\begin{align}
36 
  &= R^2 + r^2 + 2Rr \frac{\sqrt{3}-1}{2\sqrt{2}} \\
  &= R^2 + r^2 + \frac{\sqrt{3}-1}{\sqrt{2}} Rr \label{eq:2}
\end{align}
となる．\cref{eq:1,eq:2}を辺々引くと
\begin{align}
   & 20 = \left( \frac{\sqrt{6}-\sqrt{2}}{2} + 1 \right) Rr \\
\therefore 
   & Rr = \frac{40}{\sqrt{6}-\sqrt{2}+2} \label{eq:3}
\end{align}
である．

また\cref{eq:1}を変形して
\begin{align}
  R^2 + r^2 = 16 + Rr \label{eq:4}
\end{align}
である．

もとめる長さ $l\ (>0)$ として4時の状態に余弦定理を用いると
\begin{align}
l^2 = R^2 + r^2 - 2Rr \cos \frac{2\pi}{3} = R^2 + r^2 + Rr 
\end{align}
だから，\cref{eq:3,eq:4}を代入して
\begin{align}
l^2 &= (16 + Rr) + Rr \quad (\because \text{\cref{eq:4}}) \\
&= 16 + 2Rr \\
&= 16 + \frac{80}{\sqrt{6}-\sqrt{2}+2} \quad (\because \text{\cref{eq:3}}) \\
&= 16 + 20(\sqrt{3}+1-\sqrt{2}) \\
&= 20(\sqrt{3}-\sqrt{2}) + 36 \label{eq:5}
\end{align}
を得るから，この少数第一位までの評価を求めれば良い．

以下$(6.5)^2 < l^2 < (6.6)^2$ であることを示す．
\begin{align*}
(1.73)^2 < 3 < (1.733)^2 \quad &\therefore 1.73 < \sqrt{3} < 1.733 \\
(1.41)^2 < 2 < (1.42)^2 \quad &\therefore 1.41 < \sqrt{2} < 1.42
\end{align*}
だから，\cref{eq:5}に代入して
\begin{align}
 & 20(1.733 - 1.42) + 36 < l^2 < 20(1.74 - 1.41) + 36 \\
\therefore 
 &42.26 < l^2 < 42.6
\end{align}
一方で$(6.5)^2 = 42.25$, $(6.6)^2 = 43.56$ だから，
\begin{align}
 (6.5)^2 < l^2 < (6.6)^2
\end{align}
よって示された．$l > 0$ より，
\begin{align}
6.5 < l < 6.6 
\end{align}
つまり求める近似値は
\begin{align}
l \fallingdotseq 6.5 
\end{align}
である．

\newpage

\end{document}
